\documentclass[a4paper,11pt]{book}
\usepackage{listings}
\usepackage[utf8]{inputenc}
\usepackage{titlesec}
\usepackage{fancyhdr}
\usepackage[spanish,es-tabla]{babel}
\usepackage[hidelinks]{hyperref}
\usepackage{xcolor}
\usepackage{pdfpages}

% Información reutilizable
\newcommand{\asunto}{Trabajo de Fin de Máster}
\newcommand{\titulo}{Datos abiertos en la universidad}
\newcommand{\tituloEng}{Open data at University}
\newcommand{\master}{Máster Universitario en Ingeniería Informática}
\newcommand{\autor}{Germán Martínez Maldonado}
\newcommand{\email}{germaaan@gmail.com}
\newcommand{\tutor}{Juan Julián Merelo Guervós}
\newcommand{\escuela}{Escuela Técnica Superior de Ingenierías Informática y de Telecomunicación}
\newcommand{\universidad}{Universidad de Granada}
\newcommand{\ciudad}{Granada}
\newcommand{\vers}{Versión 0.1}

% Información archivo
\hypersetup{
	pdfauthor = {\autor\ (\email)},
	pdftitle = {\titulo},
	pdfsubject = {\asunto},
	pdfkeywords = {datos abiertos, web semántica, CKAN, RDF, SPARQL, software libre},
	pdfcreator = {LaTeX con el paquete texlive},
	pdfproducer = {pdflatex}
}

% Estilo de cabeceras
\pagestyle{fancy}
\fancyhf{}
\fancyhead[LO]{\leftmark}
\fancyhead[RE]{\rightmark}
\fancyhead[RO,LE]{\textbf{\thepage}}
\setlength{\headheight}{1.5\headheight}

% Redefinición de comandos
\renewcommand{\lstlistingname}{Fragmento de código}
\renewcommand{\lstlistlistingname}{Índice de fragmentos de código}
\renewcommand{\chaptermark}[1]{\markboth{\textbf{#1}}{}}
\renewcommand{\sectionmark}[1]{\markright{\textbf{\thesection. #1}}}

% Definición de colores
\definecolor{gray97}{gray}{.97}
\definecolor{gray75}{gray}{.75}
\definecolor{gray45}{gray}{.45}
\definecolor{gray30}{gray}{.94}
\definecolor{lightgray}{rgb}{.9,.9,.9}
\definecolor{darkgray}{rgb}{.4,.4,.4}
\definecolor{purple}{rgb}{0.65, 0.12, 0.82}
\definecolor{background}{HTML}{EEEEEE}
\definecolor{delim}{RGB}{20,105,176}
\colorlet{punct}{red!60!black}
\colorlet{numb}{magenta!60!black}

% Listados
\lstset{
	aboveskip=0.5cm,
	backgroundcolor=\color{gray97},
	basicstyle=\scriptsize\ttfamily,
	breaklines=true,
	commentstyle=\color{gray45},
	frame=Ltb,
	framerule=0.5pt,
	framesep=0pt,
	framexbottommargin=3pt,
	framexleftmargin=0.1cm,
	framextopmargin=3pt,
	keywordstyle=\bfseries,
	numberfirstline = false,
	numbers=left,
	numbersep=6pt,
	numberstyle=\tiny,
	rulesep=.4pt,
	rulesepcolor=\color{black},
	showstringspaces = false,
	stringstyle=\ttfamily,
	literate={á}{{\'a}}1
	         {é}{{\'e}}1
	         {í}{{\'i}}1
	         {ó}{{\'o}}1
	         {ú}{{\'u}}1
	         {ñ}{{\~n}}1
}
 
% Minimizar fragmentado de listados
\lstnewenvironment{listing}[1][]
	{\lstset{#1}\pagebreak[0]}{\pagebreak[0]}

% Para que las páginas en blanco no tengan cabecera
\makeatletter
\def\clearpage{%
  \ifvmode
    \ifnum \@dbltopnum =\m@ne
      \ifdim \pagetotal <\topskip
        \hbox{}
      \fi
    \fi
  \fi
  \newpage
  \thispagestyle{empty}
  \write\m@ne{}
  \vbox{}
  \penalty -\@Mi
}
\makeatother

\begin{document}
\begin{titlepage}

\newlength{\centeroffset}
\setlength{\centeroffset}{-0.5\oddsidemargin}
\addtolength{\centeroffset}{0.5\evensidemargin}

\noindent\hspace*{\centeroffset}\begin{minipage}{\textwidth}

\centering
\includegraphics[width=0.9\textwidth]{../images/logo_ugr.png}\\[1.4cm]

\textsc{\Large\asunto\\[0.2cm]}
\textsc{\master}\\[1cm]

{\Huge\bfseries \titulo\\}
\noindent\rule[-1ex]{\textwidth}{3pt}\\[3.5ex]
\end{minipage}

\vspace{2cm}
\noindent\hspace*{\centeroffset}\begin{minipage}{\textwidth}
\centering

\textbf{Autor}\\ {\autor}\\[2.5ex]
\textbf{Tutor}\\ {\tutor}\\[2cm]
\includegraphics[width=0.3\textwidth]{../images/logo_etsiit.png}\\[0.1cm]
\textsc{\escuela}\\
\textsc{---}\\
\ciudad, \today\\
\includegraphics[width=0.3\textwidth]{../images/CC-SA-logo.png}
\end{minipage}
\end{titlepage}
\frontmatter
\begin{center}
{\LARGE\bfseries\titulo}\\
\end{center}
\begin{center}
\autor\
\end{center}

\section*{Resumen}

\newpage
\begin{center}
{\LARGE\bfseries\tituloEng}\\
\end{center}
\begin{center}
\autor\
\end{center}

\section*{Abstract}

\newpage
\thispagestyle{empty}
\
\vspace{3cm}

\noindent\rule[-1ex]{\textwidth}{2pt}\\[4.5ex]

Yo, \textbf{\autor}, alumno de la titulación \textbf{\master} de la \textbf{\escuela\ de la \universidad}, autorizo la ubicación de la siguiente copia de mi Trabajo Fin de Máster (\textit{\titulo}) en la biblioteca del centro para que pueda ser consultada por las personas que lo deseen.

\bigskip
Además, este mismo trabajo es realizado bajo licencia \textbf{Creative Commons Attribution-ShareAlike 4.0} (\url{https://creativecommons.org/licenses/by-sa/4.0/}), dando permiso para copiarlo y redistribuirlo en cualquier medio o formato, también de adaptarlo de la forma que se quiera, pero todo esto siempre y cuando se reconozca la autoría y se distribuya con la misma licencia que el trabajo original. El documento en formato {\tt LaTeX} se puede encontrar en el siguiente repositorio de {\tt GitHub}: \url{https://github.com/germaaan/ProyectoDatosAbiertos}.

\vspace{4cm}

\noindent Fdo: \autor

\vspace{2cm}

\begin{flushright}
\ciudad, a \today
\end{flushright}

\newpage
\thispagestyle{empty}
\
\vspace{3cm}

\noindent\rule[-1ex]{\textwidth}{2pt}\\[4.5ex]

D. \textbf{\tutor}, profesor del \textbf{Departamento de Arquitectura y Tecnología de Computadores} de la \textbf{\universidad}.

\vspace{0.5cm}

\vspace{0.5cm}

\textbf{Informa:}

\vspace{0.5cm}

Que el presente trabajo, titulado \textit{\textbf{\titulo}}, ha sido realizado bajo su supervisión por \textbf{\autor}, y 
autoriza la defensa de dicho trabajo ante el tribunal que corresponda.

\vspace{0.5cm}

Y para que conste, expide y firma el presente informe en \ciudad\ a \today.

\vspace{1cm}

\textbf{El tutor:}

\vspace{5cm}

\noindent \textbf{\tutor}

\chapter*{Agradecimientos}
\thispagestyle{empty}

\vspace{1cm}

Agradecimientos.

\begingroup
\let\cleardoublepage\clearpage
  \tableofcontents
  \listoffigures
  \listoftables
  \lstlistoflistings
\endgroup

\newpage
\thispagestyle{empty}
\
\mainmatter
\chapter{Introducción}

Introducción.
\chapter{Objetivos}

Objetivos.
\chapter{Análisis y Diseño}

\section{Conceptos del sistema}
\chapter{Implementación}

Implementación.
\chapter{Pruebas}

Pruebas.
\chapter{Conclusiones y trabajos futuros}

Conclusiones y trabajos futuros.

\backmatter
\begingroup
	\setlength\parindent{0pt}
	\chapter{Glosario de términos}

Glosario
\endgroup

\addcontentsline{toc}{chapter}{Anexo}
\chapter*{Anexo}
%\input{back/manual_usuario}

\newpage
\begin{thebibliography}{99}
	\addcontentsline{toc}{chapter}{Bibliografía}

\subsubsection*{Referencias bibliográficas consultadas:}

\bibitem{SWB} ``Semantic Web roadmap''. Tim Berners-Lee. 14/10/1998. \url{https://www.w3.org/DesignIssues/Semantic.html}

\bibitem{OLSW} ``Ontology Learning for the Semantic Web. IEEE Intelligent Systems. Volume 16, Issue 2, March/April 2001, Pages 72-79''. Alexander Maedche, Steffen Staab. Marzo/Abril 2001.

\bibitem{ASW} ``Agents and the Semantic Web. IEEE Intelligent Systems. Volume 16, Issue 2, March/April 2001, Pages 30-37''. Marzo/Abril 2001.

\bibitem{SWR} ``Semantic Web: The roles of XML and RDF. IEEE Internet Computing. Volume 4, Issue 5, September 2000, Pages 63-74''. Stefan Decker, Sergey I. Melnik, Frank Van Harmelen, Dieter Fensel, Michel C.A. Klein, Jeen Broekstra, Michael A. Erdmann, Ian Horrocks. Marzo 2003.

\bibitem{OMSA} ``Ontology mapping: The state of the art. Knowledge Engineering Review. Volume 18, Issue 1, March 2003, Pages 1-31''. Yannis Kalfoglou, Marco Schorlemmer. Marzo 2003.

\bibitem{LDDI} ``Linked Data - Design Issues''. Tim Berners-Lee. 18/06/2009. \url{https://www.w3.org/DesignIssues/LinkedData.html}

\bibitem{PSW} ``Programming the Semantic Web''. Toby Segaran, Colin Evans, Jamie Taylor. Julio 2009.

\bibitem{MMCC} ``Metodologías y métodos para la construcción de ontologías. Scientia Et Technica. Nº 50, Abril de 2012. Páginas 133-140''. Abril 2011.

\bibitem{LSQU} ``Learning SPARQL: Querying and Updating with SPARQL 1.1''. Bob DuCharme. Julio 2013.

\bigskip
\subsubsection*{Recursos web consultados}

\bibitem{OWOL} ``OWL Web Ontology Language''. W3C, 10/02/2004. \url{https://www.w3.org/TR/owl-features/}

\bibitem{RDFS} ``SPARQL 1.1 Query Language''. W3C, 21/03/2013. \url{https://www.w3.org/TR/sparql11-query/}

\bibitem{RDFCAS} ``RDF 1.1 Concepts and Abstract Syntax''. W3C, 25/02/2014. \url{https://www.w3.org/TR/2014/REC-rdf11-concepts-20140225/}

\bibitem{RDFS} ``RDF Schema 1.1''. W3C, 25/02/2014. \url{https://www.w3.org/TR/2014/REC-rdf-schema-20140225/}
\end{thebibliography}

\newpage \
\thispagestyle{empty}
\end{document}